\chapter*{Abstract}
\addcontentsline{toc}{chapter}{Abstract}
\markboth{Abstract}{Abstract}

Current 3D scanning solutions available for indoor locations either need costly LiDAR hardware or complicated step-up procedures that are impractical for construction site personnel, building administrators, and small design companies. The project \textbf{BuildScan~3D} is closing the gap by offering a total system capable of converting a footage captured by a usual Android smartphone (such as the realme phone, which has no LiDAR unit or any special equipment) into a metric-scale 3D Gaussian Splatting (3DGS) model, i.e. the one that can be taken for real-world measurements.

The proposed system has five steps: (1) an individual Android app that uses ARCore to record video at the rate of 30\,FPS together with camera poses, depth maps, and image quality feedback; (2) a pose-prior stage that relates ARCore OpenGL coordinates to COLMAP/Nerfstudio ones; (3) a cloud-GPU reconstruction stage for which COLMAP Structure-from-Motion with Sim3 alignment of ARCore coordinates has been applied, and training has been performed with the use of Nerfstudio; (4) a structural layer that makes it possible to extract walls, corners, and lengths of walls and doors to create floor plans and measure objects (prototype built); (5) a planned web viewer for measurements and inspecting the model.

At the mid-point evaluation, Iteration~1 stands at \emph{core pipeline complete, refinement ongoing}: the Android capture application processes all the necessary sensor data with live quality and coverage guidance, the FastAPI processing server receives uploads and sends work to RunPod GPU pods as needed, and the top reconstruction (version six) achieved 26.52\,dB PSNR/0.8886 SSIM on forty-two held-out frames; tape measure calculations show an average measurement error of 4.8\,mm (0.19\%) and a maximum error of 0.29\% which is far better than the five percent accuracy requirement. The most important technical issues that have been solved already are bridging the tracking gap with IMU, lens self-calibration using COLMAP, and discovering and fixing the problems with irregular Gaussian splat geometry. The open items remaining include refining the 3DGS, creating 3D mesh, completing the structural layer and creating a web workspace (including the floor maps and 1:1 measurements), which should be followed by full-scale field validation.
